\section{Preliminaries}

First some notation. Let $\N = \{1,2,\ldots\}$. For $k \in \Z_{\geq0}$, define $[k] := \{1,\ldots,k\}$. Let $I$ be the identity matrix of any size and of any vector space, which will be clear from context. For a graph $G$, let $V(G)$ denote its vertex set and $E(G)$ its edge set. For any vector $v$ with $n$ entries, let $v_i$ be its $i$th entry for all $i \in [n]$. 

\subsection{Qubits}

\begin{definition}
A single \term{qubit} $\ket\phi$ is a state
\begin{align}
    \ket\phi = \alpha\vect{1}{0} + \beta\vect{0}{1} \in \C^2,
\end{align}
where $|\alpha|^2 + |\beta|^2 = 1$ and $\alpha, \beta \in \C$. We say $\ket\phi$ is in a \term{superposition} between $\svect{1}{0}$ and $\svect{0}{1}$. Thus a single qubit is a length 2 vector in the Hilbert space $\C^2$ with a $\ell_2$ norm of 1.
\end{definition}

\begin{definition}
Define kets $\ket{0} := \svect{1}{0}$ and $\ket{1} := \svect{0}{1}$, both in $\C^2$. For any $\ket\phi$, let bra $\bra\phi := \ket\phi^\dagger = (\ket\phi^*)^T$ be its conjugate transpose. Thus the braket denotes the inner product:
\begin{align}
    \braket{\phi}{\phi}
    := \bra{\phi} \cdot \ket{\phi}.
\end{align}
Also, the $\ell_2$ norm or length $|\ket{\phi}| := \sqrt{\braket{\phi}{\phi}}$. In general, for any $\ket\phi$ and $\ket\psi$,
\begin{align}
    \braket{\phi}{\psi}
    := \bra{\phi} \cdot \ket{\psi}
    = \sum_{i=1}^{2} \phi_i^* \psi_i.
\end{align}
\end{definition}

\begin{definition}
Let $A \in \C^{m \times n}$ and $B \in \C^{o \times p}$. For our purposes, it suffices to define the \term{tensor product}
\begin{align}
    A \otimes B := \begin{pmatrix}
        a_{1,1}B & \cdots & a_{1,n} \\
        \vdots & \ddots & \vdots \\
        a_{n,1}B & \cdots & a_{n,m}B
    \end{pmatrix}
\end{align}
where $a_{i,j}$ denotes the $i$th row and $j$th column of $A$. If $x = \overline{x_1 x_2 \cdots x_d} \in \{0,1\}^d$,
\begin{align}
    \ket{x} := \bigotimes_{i=1}^d \ket{x_i}.
\end{align}
We may replace $x$ by its integer evaluation as a binary number. Also, we may write $\ket{\phi} \ket\psi = \ket\phi \otimes \ket\psi$, so adjacent kets imply $\otimes$. Also, let $A^{\otimes n} := \bigotimes_{i=1}^n A$. This gives us the definition of a $d$-qubit state
\begin{align}
    \ket\varphi = \sum_{x \in \{0,1\}^d} \alpha_x \ket{x} \in \C^{2^d},
\end{align}
where $\alpha_x \in \C$ and $\sum_{x \in \{0,1\}^d} \alpha_x \ket{x} = 1$. As the number of dimensions increases exponentially with regards to the number of qubits, braket notation will become increasingly concise compared to traditional vector notation.
\end{definition}

\begin{definition} Vectors $\ket{v_1}, \ldots, \ket{v_n} \in \C^n$ are an orthonormal basis if
\begin{enumerate}
    \item $|\ket{v_i}| = 1$ for every $i \in [n]$ and
    \item $\ket{v_i} \bot \ket{v_j}$ (equivalently, $\braket{v_i}{v_j} = 0$) for every $i,j \in [n]$ with $i \neq j$.
\end{enumerate}
Henceforth, all bases are orthonormal unless explicitly said otherwise.
\end{definition}

\begin{definition} 
The \term{Born Rule} defines how quantum measurement occurs: A $n$-qubit state $\ket\phi$ can be measured under any basis $\ket{v_i} \in \C^{2^n}$ with $i \in [2^n]$. When measured, the output $i$ (interpreted as classical bits) occurs with probability $| \braket{\phi}{v_i} |^2$. After seeing the outcome $v_i$, the state collapses to $\ket{v_i}$. This process is irreversible.
\end{definition}

\begin{definition}
One can also view measurement as projection. This is useful for \term{partial measurement}, where only some qubits in a state are measured. Let the $n$ qubit space $\C^{2^n} = S_1 \otimes S_2$ where $S_1 \cong (\C^2)^k$. Let $\{\ket{s_i}\}_{i \in [k]}$ be a basis for $S_1$. Define the \term{projection matrix} on $S_1$ as
\begin{align}
    P := \sum_{j=1}^k \ket{s_j} \bra{s_j}, \quad 
    \text{such that} \quad
    P = P^\dagger, \quad
    P^2 = P.
\end{align}
Define the \term{partial measurement} on $S_1$ by $M = P \otimes I \in \C^{2^n \times 2^n}$. For a state $\ket\phi \in \C^{2^n}$, 
\begin{align}
    \Pr(\text{Measuring }i\text{ from}\ket\phi)
    = \bra\phi M \ket\phi 
    = \| M \ket\phi \|^2.
\end{align}
On measuring outcome $i$, the post-measurement state becomes
\begin{align}
    \ket{\phi_i}
    = \frac{M \ket\phi}{\sqrt{\bra\phi M \ket\phi}}.
\end{align}
In other words, $\ket{\phi_i}$ is the state that minimizes $\|\ket{\phi} - \ket{\phi_{j}}\|$ over all $\ket{\phi_j} \in S_1$. Partial measurements technically give rise to mixed states, which cannot be represented as a single ket, but we defer its discussion to a later section.
\end{definition}

\begin{example}
Consider the states
\begin{align}
    \ket{+} := \frac{1}{\sqrt2}\ket0 + \frac{1}{\sqrt2}\ket1, \quad
    \ket{-} := \frac{1}{\sqrt2}\ket0 - \frac{1}{\sqrt2}\ket1.
\end{align}
Measuring either in the \term{standard basis} $\{\ket0, \ket1\}$ gives $\ket0$ with $0.5$ probability and $\ket1$ with $0.5$ probability since
\begin{align*}
    |\braket{+}{0}|^2 
    = \left| \frac{1}{\sqrt2} \cdot 1 + \frac{1}{\sqrt2} \cdot 0 \right|^2
    = \frac{1}{2}
\end{align*}
and similarly for $\ket1$. The squaring makes $\ket{-}$ proceed similarly as well.
\end{example}

\begin{definition}
A two qubit state $\ket\phi$ is \term{entangled} if it cannot be written as the tensor product of single qubits. That is, $\nexists \ket\psi, \ket\varphi \in \C^2$ such that $\ket\phi = \ket\psi \otimes \ket\varphi$. If a state is not entangled, it is called \term{separable}.
\end{definition}

\subsection{Quantum Gates}

\begin{definition}
A single qubit quantum gate $U \in \C^{2 \times 2}$ satisfies one of:
\begin{enumerate}
    \item $U^\dagger U = UU^\dagger = I$;
    \item equivalently, $U$ is a unitary matrix;
    \item equivalently, $U = \ket{\phi_0}\bra0 + \ket{\phi_1}\bra1$, where $\{\ket{\phi_0}, \ket{\phi_1}\}$ form an orthonormal basis.
\end{enumerate}
Common single qubit gates are given in Figure \ref{fig:single-qubit-gates}.
\end{definition}

\begin{figure}[H]
\centering
\renewcommand{\arraystretch}{1.1}
\begin{tabular}{c c c c}

\toprule
\textbf{Name} & \textbf{Circuit} & \textbf{Matrix} & \textbf{Bra-Ket} \\
\midrule

$X$ &
$\begin{quantikz}
\lstick{} & \targ{} & \qw
\end{quantikz}$ &
$\begin{pmatrix}
0 & 1 \\
1 & 0
\end{pmatrix}$ &
$\ket{1}\bra{0} + \ket{0}\bra{1}$ \\

$Y$ &
$\begin{quantikz}
\lstick{} & \gate{Y} & \qw
\end{quantikz}$ &
$\begin{pmatrix}
0 & -i \\
i & 0
\end{pmatrix}$ &
$i\ket{1}\bra{0} - i\ket{0}\bra{1}$ \\

$Z$ &
$\begin{quantikz}
\lstick{} & \gate{Z} & \qw
\end{quantikz}$ &
$\begin{pmatrix}
1 & 0 \\
0 & -1
\end{pmatrix}$ &
$\ket{0}\bra{0} - \ket{1}\bra{1}$ \\

$H$ &
$\begin{quantikz}
\lstick{} & \gate{H} & \qw
\end{quantikz}$ &
$\displaystyle\frac{1}{\sqrt{2}}
\begin{pmatrix}
1 & 1 \\
1 & -1
\end{pmatrix}$ &
$\ket{+}\bra{0} + \ket{-}\bra{1}$ \\

$S$ &
$\begin{quantikz}
\lstick{} & \gate{S} & \qw
\end{quantikz}$ &
$\begin{pmatrix}
1 & 0 \\
0 & i
\end{pmatrix}$ &
$\ket{0}\bra{0} + i\ket{1}\bra{1}$ \\

$T$ &
$\begin{quantikz}
\lstick{} & \gate{T} & \qw
\end{quantikz}$ &
$\begin{pmatrix}
1 & 0 \\
0 & e^{i\pi/4}
\end{pmatrix}$ &
$\ket{0}\bra{0} + e^{i\pi/4}\ket{1}\bra{1}$ \\

$R_x(\theta)$ &
$\begin{quantikz}
\lstick{} & \gate{R_x(\theta)} & \qw
\end{quantikz}$ &
$\begin{pmatrix}
\cos(\nicefrac{\theta}{2}) & -i\sin(\nicefrac{\theta}{2}) \\
-i\sin(\nicefrac{\theta}{2}) & \cos(\nicefrac{\theta}{2})
\end{pmatrix}$ &
$\begin{aligned}
\cos(\nicefrac{\theta}{2})\ket{0}\bra{0} 
&- i\sin(\nicefrac{\theta}{2})\ket{1}\bra{0} \\[-0.5ex]
- i\sin(\nicefrac{\theta}{2})\ket{0}\bra{1} 
&+ \cos(\nicefrac{\theta}{2})\ket{1}\bra{1}
\end{aligned}$ \\

$R_y(\theta)$ &
$\begin{quantikz}
\lstick{} & \gate{R_y(\theta)} & \qw
\end{quantikz}$ &
$\begin{pmatrix}
\cos(\nicefrac{\theta}{2}) & -\sin(\nicefrac{\theta}{2}) \\
\sin(\nicefrac{\theta}{2}) & \cos(\nicefrac{\theta}{2})
\end{pmatrix}$ &
$\begin{aligned}
\cos(\nicefrac{\theta}{2})\ket{0}\bra{0} 
&+ \sin(\nicefrac{\theta}{2})\ket{1}\bra{0} \\[-0.5ex]
- \sin(\nicefrac{\theta}{2})\ket{0}\bra{1}
&+ \cos(\nicefrac{\theta}{2})\ket{1}\bra{1}
\end{aligned}$ \\

$R_z(\theta)$ &
$\begin{quantikz}
\lstick{} & \gate{R_z(\theta)} & \qw
\end{quantikz}$ &
$\begin{pmatrix}
e^{-i\nicefrac{\theta}{2}} & 0 \\
0 & e^{i\nicefrac{\theta}{2}}
\end{pmatrix}$ &
$\begin{aligned}
e^{-i\nicefrac{\theta}{2}}\ket{0}\bra{0}
+ e^{i\nicefrac{\theta}{2}}\ket{1}\bra{1}
\end{aligned}$ \\

\bottomrule

\end{tabular}

\caption{Common single qubit quantum gates}
\label{fig:single-qubit-gates}
\end{figure}
\renewcommand{\arraystretch}{1}


\begin{fact} 
\label{fact:tensor-properties}
For matrices $A,B,C,D$, scalar $\lambda$, and vectors $\ket\phi$, $\ket\psi$, $\otimes$ satisfies:
\begin{enumerate}
    \item Bi-linearity: $(A \otimes B)(C \otimes D) = (AC) \otimes (BD)$, so $(A \otimes B)(\ket\phi \otimes \ket\psi) = (A\ket\phi) \otimes (B\ket\psi)$.
    \item Scalars commute: $\lambda(A \otimes B) = (\lambda A) \otimes B = A \otimes (\lambda B)$.
    \item Associativity: $(A \otimes B) \otimes C = A \otimes (B \otimes C)$.
    \item Unitary-preserving: If $A,B$ are unitary matrices, $A \otimes B$ is unitary.
\end{enumerate}
\end{fact}

\begin{definition} The general $1$-qubit unitary is given by
\begin{align}
    U(\theta,\iota,\kappa)
    :=
    \begin{pmatrix}
    \cos\!\left(\frac{\theta}{2}\right)
    & - e^{i\kappa}\sin\!\left(\frac{\theta}{2}\right) \\
    e^{i\iota}\sin\!\left(\frac{\theta}{2}\right)
    & e^{i(\iota+\kappa)}\cos\!\left(\frac{\theta}{2}\right)
    \end{pmatrix}.
\end{align}
\end{definition}

\begin{fact}
Any 1-qubit gate can be written as $U(\theta,\iota,\kappa)$ for some $\theta,\iota,\kappa \in [0,2\pi]$ [cite:need].
\end{fact}

\begin{definition}
An \term{$n$-qubit gate} is a unitary matrix in $\C^{2^n \times 2^n}$. If gates $U_1, \ldots, U_k$ satisfy $U = \bigotimes_{i=1}^k U_i \in \C^{2^n}$, by Fact \ref{fact:tensor-properties}.4 $U$ is an $n$-qubit gate. Special named multi-qubit gates are given in Figure \ref{fig:multi-qubit-gates}.
\end{definition}

\begin{figure}[H]
\centering
\renewcommand{\arraystretch}{1.1}
\begin{tabular}{c c c}

\toprule
\textbf{Name} & \textbf{Circuit} & \textbf{Bra-Ket} \\
\midrule

CNOT &
$\begin{quantikz}
\lstick{} & \ctrl{1} & \qw \\
\lstick{} & \targ{} & \qw
\end{quantikz}$ &
$\ket{0}\bra{0}\otimes I + \ket{1}\bra{1}\otimes X$ \\

CCNOT or Toffoli &
$\begin{quantikz}
\lstick{} & \ctrl{1} & \qw \\
\lstick{} & \ctrl{1} & \qw \\
\lstick{} & \targ{} & \qw
\end{quantikz}$ &
$\displaystyle \ket{11}\bra{11}\otimes X + \sum_{\substack{(a,b)\neq(1,1)\\a,b \in \{0,1\}}} \ket{ab}\bra{ab}\otimes I$ \\

SWAP &
$\begin{quantikz}
\lstick{} & \swap{1} & \qw \\
\lstick{} & \targX{} & \qw
\end{quantikz}$ &
$\displaystyle\sum_{i,j\in\{0,1\}} \ket{ij}\bra{ji}$ \\

C-$Z$ &
$\begin{quantikz}
\lstick{} & \ctrl{1} & \qw \\
\lstick{} & \ctrl{-1} & \qw
\end{quantikz}$ &
$\ket{0}\bra{0}\otimes I + \ket{1}\bra{1}\otimes Z$ \\

\bottomrule

\end{tabular}

\caption{Common multi-qubit quantum gates}
\label{fig:multi-qubit-gates}
\end{figure}
\renewcommand{\arraystretch}{1}


\begin{definition}
Many of the gates in Figure \ref{fig:multi-qubit-gates} are \term{controlled} gates. Informally, controlled gates, such as $\begin{quantikz}&\targ{0}&\end{quantikz}$ only trigger if all of the controls $\begin{quantikz}&\ctrl{0}&\end{quantikz}$ connected to it are $\ket{1}$. Formally, for any $n$-qubit gate $U$, the single controlled gate, denoted $\cl U$, is defined
\begin{align}
    \cl U := \ket0 \bra0 \otimes I + \ket1 \bra1 \otimes U.
\end{align}
This definition generalizes for multiple controls. Thus $\CNOT$ is just $\cl X$ and Toffoli is $\text{CC-}X$. We may suppress the - if the meaning is clear. $\cl Z$ is special as it only triggers if both qubits are $\ket1$. Note that $I = \ket0 \bra0 + \ket1 \bra1$ so $Z = I - 2 \ket1 \bra1$. Then
\begin{align*}
    \cl Z 
    = \ket0 \bra0 \otimes I + \ket1 \bra1 \otimes (I - 2 \ket1 \bra1)
    = I \otimes I - 2 \ket{11} \bra{11}.
\end{align*}
Thus, it can be interpreted that both qubits of the $\cl Z$ are controls.
\end{definition}

\subsection{Universality}

\begin{definition}
For every $n \in \N$, Let $\SU(n)$ denote the special unitary group under matrix multiplication consisting of $n \times n$ unitary matrices with determinant 1. It is a normal subgroup of the unitary group $U(n)$, consisting of $n \times n$ unitary matrices.
\end{definition}

\begin{definition}
\label{def:universality}
A set of gates $\cS$ is (efficiently) \term{universal} if any $U \in \SU(2^n)$, for any $n \in \N$, that can be generated (in $\Poly(n)$ gates) by an arbitrary gate set $\cS'$, can also be generated by the gate set $\cS$ (in $\Poly(n)$ gates) \cite{MBQC-universal}. We say a set of gates $\cS$ is $d$-qubit universal if the above holds for $U \in \SU(2^d)$.
\end{definition}

\begin{remark}
In this thesis we will only focus on universality, not efficient universality. Thus, it may take exponential or more gates to simulate certain unitaries $U$. Furthermore, in practice, the application of $U$ up to some measure and error $\epsilon$ suffices. This is called \term{approximate universality}. This thesis will not focus on approximate universality either.
\end{remark}

\begin{remark}
Definition \ref{def:universality} is only concerned with implementing arbitrary unitaries, not necessarily computation. That is, it does not specify whether the inputs and ouputs are classical or quantum in nature. This matters because $\ket\phi = \ket0$ and $\ket\psi = -\ket0$ have the same measurement distribution and classical output, but $\ket\phi \neq \ket\psi$. This difference is called \term{global phase}, i.e. a scalar $e^{i\theta}$ multiplied to an entire state.
\end{remark}

\begin{definition}
We use the following types of \term{universal computation} from \cite{MBQC-universal}:
\begin{itemize}
    \item CC-universal: For every $U \in \SU(n)$, the model can reproduce the distribution of local measurements performed on $U \{\ket{0^n}\}$.
    \item QC-universal: Given an arbitrary $n$-qubit state $\ket\phi$, for every $U \in \SU(2^n)$, the model can reproduce the distribution of local measurements performed on $U \ket\phi$.
    \item CQ-universal: For every $U \in SU(2^n)$, the model can produce the quantum state $\ket\phi = U\ket{0^n}$.
    \item QQ-universal: Given an arbitrary $n$-qubit state $\ket\phi$, for every $U \in \SU(2^n)$, the model can produce the quantum state $U \ket\phi$.
\end{itemize}
\end{definition}

\begin{fact}
Any 1-qubit universal gate set and a 2 qubit entangling gate is universal [cite:need].
\end{fact}

\begin{fact}
Any universal gate set must include at least one multi-qubit entangling gate, as it is impossible to entangle using only single-qubit gates.
\end{fact}

\begin{example}
The \term{Clifford gates} consist of Hadamard $H$, Phase $S$, and $\CNOT$. The Gottesman-Knill theorem \cite{Gottesman-Knill} shows this gate set is not universal and is classically efficiently simulable.
\end{example}

\subsection{Measurement-Based Quantum Computation}

\begin{definition}
Measurement-based quantum computation (MBQC) roughly proceeds as follows:
\begin{enumerate}
    \item A classical input specifies the data and program.
    \item A resource state $\ket{\phi_k}$ is chosen from some set of states $\Phi = \{ \ket{\phi_1}, \ket{\phi_2}, \ldots \}$.
    \item Local operations (single qubit unitaries), local measurements, and classical communication are used to transform this resource state.
    \item By the end, we have a state $\ket\psi \otimes \ket{\junk}$. This $\ket\psi$ or its measurement is the desired result.
\end{enumerate}
The rest of this section formalizes these notions.
\end{definition}

\begin{definition}
A resource state is a quantum state $\ket\phi$ that can be used as the initial state in MBQC. A \term{graph state} $\ket{G}$ is a particular type of resource state associated with an undirected graph $G = (V,E)$ such that:
\begin{itemize}
    \item For every $v \in V(G)$, create a new qubit in the $\ket+$ state using $H$.
    \item For every $e = (u,v) \in E(G)$, apply $\cl Z$ on qubits associated with $u$ and $v$.
\end{itemize}
\end{definition}

\begin{fact}
There is a unique construction for every graph state, up to reordering of $\cl Z$. This is because $\cl Z$s commute with $I$ and $\cl Z$.
\begin{proof}
The $\cl Z$ gate is a diagonal matrix, and so is $I$. The tensor product of diagonal matrices is a diagonal matrix. Furthermore, diagonal matrices commute. Thus the order of adding edges, or equivalently applying $\cl Z$s, does not matter.
\end{proof}
\end{fact}

\begin{fact}
Let $G$ be a graph with maximum degree $\Delta$. Then $\ket{G}$ can be constructed in depth $\Delta + 1$.
\begin{proof}
Vizing's theorem states that any graph with maximum degree $\Delta$ admits an edge coloring with $\Delta + 1$ colors such that no two edges that share a vertex have the same color. All edges of the same color can be implemented in the graph state in the same layer of $\cl Z$s. Thus, at most $\Delta + 1$ layers.
\end{proof}
\end{fact}

\begin{fact}
Define the Hadamard gadget as the following circuit:
\[
\begin{quantikz}
\lstick{$\ket{\phi}$} & \ctrl{1} & \gate{H} & \meter{} \\
\lstick{$\ket{+}$}    & \ctrl{-1} & \qw      & \qw
\end{quantikz}
\]
Commuting a Hadamard through the $\cl Z$ gives a $\CNOT$:
\[
\begin{aligned}
\begin{quantikz}
\lstick{$\ket{\phi}$} & \ctrl{1} & \gate{H} & \meter{} \\
\lstick{$\ket{+}$}    & \ctrl{-1} & \qw      & \qw
\end{quantikz}
\;=\;
\begin{quantikz}
\lstick{$\ket{\phi}$} & \gate{H} & \targ{}  & \meter{} \\
\lstick{$\ket{+}$}    & \qw      & \ctrl{-1}& \qw
\end{quantikz}
\;=\;
\begin{quantikz}
\lstick{$H\ket{\phi}$} & \targ{}  & \meter{} \\
\lstick{$\ket{+}$}     & \ctrl{-1}& \qw
\end{quantikz}
\end{aligned}
\]
Two $\CNOT$s cancel each other out and a $\CNOT$ has no effect on $\ket+$. Then
\[
\begin{aligned}
\begin{quantikz}
\lstick{$H\ket{\phi}$} & \targ{}  & \meter{} \\
\lstick{$\ket{+}$}     & \ctrl{-1}& \qw
\end{quantikz}
\;=\;
\begin{quantikz}
\lstick{$H\ket{\phi}$} & \ctrl{1} & \targ{}  & \ctrl{1} & \ctrl{1} & \meter{} \\
\lstick{$\ket{+}$}     & \targ{}  & \ctrl{-1}& \targ{}  & \targ{}  & \qw
\end{quantikz}
\end{aligned}
\]
The first three $\CNOT$s are equivalent to a swap gate. Thus
\[
\begin{aligned}
\begin{quantikz}
\lstick{$H\ket{\phi}$} & \ctrl{1} & \targ{}  & \ctrl{1} & \targ{}  & \meter{} \\
\lstick{$\ket{+}$}     & \targ{}  & \ctrl{-1}& \targ{}  & \ctrl{-1}& \qw
\end{quantikz}
\;=\;
\begin{quantikz}
\lstick{$\ket{+}$}     & \targ{}  & \meter{} \\
\lstick{$H\ket{\phi}$} & \ctrl{-1}& \qw
\end{quantikz}
\end{aligned}
\]
The final state before measurement is
\begin{align*}
    \CNOT (\ket+ \otimes H \ket\phi) 
    = \frac{1}{\sqrt2} \left( \ket0 \otimes H \ket\phi + \ket1 \otimes XH \ket\phi \right).
\end{align*}
In other words, measuring in the $X$ basis ($\{\ket0, \ket1\}$) applies an $X$ to $\ket\phi$ if we measure a $\ket{1}$. This happens with only $0.5$ probability, hence the need for \term{adaptive measurements}. That is, measurements whose classical outcome can dynamically impact how other qubits are measured.
\end{fact}

\begin{fact}
$R_z(\theta) Z R_z(\theta)^\dagger = Z$. That is, $R_z(\theta)$ commutes with $\cl Z$ for all $\theta$. Thus,
\[
\begin{aligned}
\begin{quantikz}
\lstick{$\ket{\phi}$} & \ctrl{1}  & \gate{R_z(\theta)} & \gate{H} & \meter{} \\
\lstick{$\ket{+}$}    & \ctrl{-1} & \qw      &          &
\end{quantikz}
\;=\;
\begin{quantikz}
\lstick{$\ket{+}$}              & \ctrl{1}  & \meter{} \\
\lstick{$HR_z{\theta}\ket\phi$} & \targ{}   & \qw
\end{quantikz}
\end{aligned}.
\]
\end{fact}

\begin{fact}
\label{fact:euler-decomposition}
[Euler decomposition] Every single-qubit unitary can be expressed (up to global phase) as the product $R_z(\theta_3)R_x(\theta_2)R_z(\theta_1)$.
\end{fact}

\begin{fact}
For all $\theta \in [0,2\pi]$,
\begin{align}
    R_x(\theta) 
    = H R_z(\theta) H.
\end{align} 
Thus by Fact \ref{fact:euler-decomposition} every single-qubit unitary can be applied to any $\ket\phi$ given we are making the right measurements.
\end{fact}

\begin{theorem}
\label{thm:mbqc-line}
[MBQC on the line is 1-qubit CQ-universal] Suppose you wish to apply gates $U_1, \ldots U_n$ to $\ket+$. There is a sequence of single-qubit measurements you can do to the line graph state of $|V| = \Theta(n)$ such that, assuming measurements are all $0$, gives $U_1, \ldots, U_n \ket+$ on the last qubit. 
\end{theorem}

The previous statement is really more of a statement on the computational power (in terms of universality) of the line graph as opposed to MBQC itself. For example, the graph with a single vertex is clearly not universal with the same MBQC gadgets. We formalize this notion below:

\begin{theorem}
\label{thm:cq-implies-qq}
[\cite{MBQC-universal} Observation 2: CQ-universal implies QQ-universal]. A resource is CQ-universal if, using the MBQC protocol (namely, local operations and measurements), we obtain CQ-universality. Every CQ-universal resource can be made QQ-universal by allowing quantum inputs to be coupled into the resource state using the same coupling method as before, i.e. $\cl Z$ gates.
\end{theorem}

\subsection{The Grid as a Resource State}

\begin{definition}
A $m \times n$ grid graph $G = (V,E)$ has
\begin{align}
    V &= \{(i,j) \mid i \in [m], j \in [n]\}, \\
    E &= \{\{(i_1,j_1),(i_2,j_2)\} \mid (i_1,j_1),(i_2,j_2) \in V, |i_2-i_1| + |j_2-j_1| = 1\}
\end{align}
In other words, $m$ rows and $n$ columns with edges between nodes of taxicab distance $1$.
\end{definition}